% ----- CHAUPAI 3 -----

\newpage

\subsection*{\deva{तृतीय चौपाई} (Third Chaupai)}
\addcontentsline{toc}{subsection}{\deva{तृतीय चौपाई} (Third Chaupai) — \deva{महाबीर बिक्रम...}}

\begin{minipage}[t]{0.55\textwidth}
\vspace{0pt}

{\large \linespread{1.5}\selectfont
\deva{महाबीर बिक्रम बजरंगी।}\\
\deva{कुमति निवार सुमति के संगी॥}
\par}

\vspace{0.5em}

{\large \linespread{1.5}\selectfont
\telu{మహాబీర బిక్రమ బజరంగీ।}\\
\telu{కుమతి నివార సుమతి కే సంగీ॥}
\par}

\vspace{0.5em}

{\large \linespread{1.3}\selectfont
\textit{mahābīra bikrama bajaraṅgī |}\\
\textit{kumati nivāra sumati ke saṅgī ||}
\par}
\end{minipage}%
\hfill
\begin{minipage}[t]{0.42\textwidth}
\vspace{0pt}
\begin{tcolorbox}[anchorbox]
\textbf{Choosing Good Company:} Physical strength (a \textit{Vajra} body) is dangerous without a clear mind. Hanuman's greatest power is that he dispels toxic, negative thoughts (\textit{Kumati}) and actively brings us into the company of good thoughts and wise people (\textit{Sumati ke Sangi}). True strength is maintaining a disciplined mind.
\vspace{0.5em}\newline When everyone else panicked, Hanuman delivered brilliant, clear-headed counsel to accept Vibhishana.\newline\null\hfill\href{https://www.valmiki.iitk.ac.in/sloka?field_kanda_tid=6&language=dv&field_sarga_value=17}{\textit{Yuddha Kanda, 17:50-68}}
\end{tcolorbox}
\end{minipage}

\vspace{0.5em}
\subsubsection*{\textbf{Visual Rhythm Map:}}
\begin{table}[h!]
\centering
\renewcommand{\arraystretch}{1.2}
\setlength{\tabcolsep}{0pt}
\begin{tabularx}{\textwidth}{|*{16}{>{\centering\arraybackslash\hsize=1.0\hsize}X|}}
\hline
\multicolumn{6}{|>{\centering\arraybackslash\hsize=6.0\hsize}X|}{\textbf{\laghu{म}\guru{हा}\guru{बी}\laghu{र}} \par (6)} &
\multicolumn{4}{>{\centering\arraybackslash\hsize=4.0\hsize}X|}{\textbf{\guru{बि}\laghu{क्र}\laghu{म}} \par (4)} &
\multicolumn{6}{>{\centering\arraybackslash\hsize=6.0\hsize}X|}{\textbf{\laghu{ब}\laghu{ज}\guru{रं}\guru{गी}} \par (6)} \\
\hline
\end{tabularx}
\vspace{-1.5em}
\end{table}

\begin{table}[h!]
\centering
\renewcommand{\arraystretch}{1.2}
\setlength{\tabcolsep}{0pt}
\begin{tabularx}{\textwidth}{|*{16}{>{\centering\arraybackslash\hsize=1.0\hsize}X|}}
\hline
\multicolumn{3}{|>{\centering\arraybackslash\hsize=3.0\hsize}X|}{\textbf{\laghu{कु}\laghu{म}\laghu{ति}} \par (3)} &
\multicolumn{4}{>{\centering\arraybackslash\hsize=4.0\hsize}X|}{\textbf{\laghu{नि}\guru{वा}\laghu{र}} \par (4)} &
\multicolumn{3}{>{\centering\arraybackslash\hsize=3.0\hsize}X|}{\textbf{\laghu{सु}\laghu{म}\laghu{ति}} \par (3)} &
\multicolumn{2}{>{\centering\arraybackslash\hsize=2.0\hsize}X|}{\textbf{\guru{के}} \par (2)} &
\multicolumn{4}{>{\centering\arraybackslash\hsize=4.0\hsize}X|}{\textbf{\guru{सं}\guru{गी}} \par (4)} \\
\hline
\end{tabularx}
\vspace{-1.5em}
\end{table}

\vspace{1em}
\textbf{ \deva{सरल-संस्कृतम्}:}\\
 \deva{हे महावीर! त्वं पराक्रमी वज्राङ्गी (वज्र-समान-अङ्गवान्) च असि।}\\
 \deva{त्वं कुमतिं निवारयसि, सुमतीनां च सङ्गी (रक्षकः) असि॥}\\

O Great Hero! You are valiant and possess a body as strong as a thunderbolt.\\
You dispel evil thoughts, and you are the companion of good thoughts.

\vspace{1em}
\newpage
\textbf{\deva{प्रतिपदार्थः} (Meaning of each word):}
\begin{table}[h!]
\centering
\renewcommand{\arraystretch}{1.2}
\begin{tabularx}{\textwidth}{@{} l l X X @{}}
\toprule
\textbf{Awadhi} & \textbf{Sanskrit} & \textbf{English} & \textbf{Grammar \& Notes} \\
\midrule
\deva{महाबीर} / \telu{మహాబీర} / \textit{mahābīra}  &  \deva{महावीरः}  &  Great Hero  & \\ 
\deva{बिक्रम} \gotchaicon / \telu{బిక్రమ} / \textit{bikrama}  &  \deva{विक्रमः}  &  Valiant/Courageous  & \\ 
\deva{बजरंगी} / \telu{బజరంగీ} / \textit{bajaraṅgī}  &  \deva{वज्राङ्गी}  &  Thunderbolt body  & Swara-Bhakti splitting \\
\deva{कुमति} / \telu{కుమతి} / \textit{kumati}  &  \deva{कुमतिम्}  &  Evil intellect  &   \\
\deva{निवार} \gotchaicon / \telu{నివార} / \textit{nivāra}  &  \deva{निवारकः}  &  Dispeller / Remover  & Meter trap: short 'i' \\
\deva{सुमति} / \telu{సుమతి} / \textit{sumati}  &  \deva{सुमतीनाम्}  &  Good intellect  &   \\
\deva{संगी} / \telu{సంగీ} / \textit{saṅgī}  &  \deva{सङ्गी / सखा}  &  Companion  &   \\
\bottomrule
\end{tabularx}
\end{table}

\begin{tcolorbox}[gotchabox]
\begin{itemize}
    \item \textbf{The Triple \deva{व} $\rightarrow$ \deva{ब} Shift:} This line is the ultimate test of the Awadhi rule! \deva{वीर} (Veera) $\rightarrow$ \deva{बीर} (Beera). \deva{विक्रम} (Vikrama) $\rightarrow$ \deva{बिक्रम} (Bikrama). 
    \item \textbf{Swara-Bhakti (\deva{वज्र $\rightarrow$ बजर}):} The harsh conjunct 'jra' in \deva{वज्र} (Vajra) is broken apart by inserting a neutralizing 'a' vowel, stretching it from Vaj-ra to Ba-ja-ra. This preserves the mouth's flow.
    \item \textbf{Meter Trap:} Keep the 'i' short in \deva{निवार} (Nivaara). If you elongate it to \deva{नीवार} (Neevaara), you add a beat and break the 16-matra math!
\end{itemize}
\end{tcolorbox}

\newpage
